\section{What the height equation requires}
\label{sec:scalar-core}

The geometric reduction leaves considerable freedom: we may choose
the rotational metric as well as its nonradial height. This section
identifies the conditions that guide those choices. A connection
form gives the scalar equation and its conservation law; its
linearization reveals the changes of type; a finite model then
separates the two resonance problems that the construction must solve.

\subsection{A curvature residual with a symmetric derivative}

Write $h'=w$, $E=1+w^2$, and introduce the local height jets
\begin{equation}\label{eq:height-jet-notation}
 k=z_r,\qquad p=z_\theta/r,\qquad
 D=E(1-p^2)-k^2,\qquad Q=k+Ep_\theta.
\end{equation}
The metric $g-dz^2$ is positive if and only if $D>0$. A useful
orthonormal coframe for it is
\[
 \eta_1=\frac{\sqrt D}{\sqrt{1-p^2}}\,dr,\qquad
 \eta_2=r\sqrt{1-p^2}\,d\theta-
             \frac{kp}{\sqrt{1-p^2}}\,dr.
\]
In this polar frame the Euclidean connection is $d\theta$.
We use $dr\,d\theta$ for all scalar adjoints; the source measure
must remain explicit when coordinates are changed later.

\begin{proposition}\label{prop:scalar-core}
For $D>0$, the connection form and the curvature residual are
\begin{align}
 \alpha&=D^{-1/2}\left(d\theta+\frac{k}{1-p^2}\,dp\right),
       \label{eq:conn-one-form}\\
 r\mathcal N_h(z)&=D^{-3/2}
 \left[Qk_r-Erp_r^2-\frac{E'}2(1-p^2+kp_\theta)\right],
       \label{eq:cart-curvature-scalar}
\end{align}
where $\mathcal N_h(z)=-\sqrt D\,K(g-dz^2)$ and
$d\alpha=r\mathcal N_h(z)\,dr\wedge d\theta$.
Holding $w$ fixed, the derivative is
\begin{equation}\label{eq:cart-true-derivative}
 rD_z\mathcal N_h(z)\eta
  =\partial_r(a\eta_r+b\eta_\theta)
       +\partial_\theta(b\eta_r+c\eta_\theta),
\end{equation}
with
\begin{equation}\label{eq:linear-coefficients}
 a=QD^{-3/2},\quad b=-Ep_rD^{-3/2},\quad
 c=r^{-1}(Ek_r-E'k/2)D^{-3/2}.
\end{equation}
Its determinant satisfies
\begin{equation}\label{eq:cart-determinant-identity}
 ac-b^2=\frac{K_0}{D^2}+ED^{-3/2}\mathcal N_h(z),
 \qquad K_0=\frac{ww'}r.
\end{equation}
\end{proposition}

\begin{proof}
Equality of mixed height derivatives gives
\begin{equation}\label{eq:height-compatibility}
 rp_r=k_\theta-p.
\end{equation}
Substitution in the two structure equations for $\eta_1,\eta_2$
gives \eqref{eq:conn-one-form}. Here is a direct check of its
curvature numerator, which also fixes all signs. Use
$D_r=E'(1-p^2)-2Epp_r-2kk_r$ and
$D_\theta=-2Epp_\theta-2kk_\theta$. The coefficient of
$dr\wedge d\theta$ in $D^{3/2}d\alpha$ is
\[
 -\tfrac12D_r+\frac{D}{1-p^2}(k_rp_\theta-k_\theta p_r)
 -\frac{k}{2(1-p^2)}(D_rp_\theta-D_\theta p_r).
\]
Replacing $k_\theta-p$ by $rp_r$ reduces it to the bracket in
\eqref{eq:cart-curvature-scalar}. This is also the Darboux numerator:
\begin{equation}\label{eq:darboux-numerator}
 Qk_r-Erp_r^2-\frac{E'}2(1-p^2+kp_\theta)=0.
\end{equation}
For example the Christoffel symbols are
$\Gamma^r_{rr}=E'/2E$, $\Gamma^r_{\theta\theta}=-r/E$,
$\Gamma^\theta_{r\theta}=1/r$, and
\begin{equation}\label{eq:radial-curvature}
 K_g=E'/(2rE^2).
\end{equation}
These formulas give the same numerator from
\eqref{eq:intro-darboux}.

For the variation, set $B_1=k/((1-p^2)\sqrt D)$. The identities
$(B_1)_k=ED^{-3/2}$ and $(B_1)_E=-k/(2D^{3/2})$ give
\begin{equation}\label{eq:conn-exact-variation}
 \delta\alpha=d(B_1\delta p)
  +(a\eta_r+b\eta_\theta)\,d\theta
  -(b\eta_r+c\eta_\theta)\,dr,
 \qquad\delta p=\eta_\theta/r.
\end{equation}
Taking the exterior derivative removes the exact term and proves
\eqref{eq:cart-true-derivative}. In particular the normal flux is
$a\eta_r+b\eta_\theta$.

Finally, let $N$ denote the bracket in
\eqref{eq:cart-curvature-scalar}. Since
$kQ+D=E(1-p^2+kp_\theta)$,
\[
 ac-b^2=\frac E{rD^3}\left(N+\frac{E'D}{2E}\right).
\]
Using $N=rD^{3/2}\mathcal N_h(z)$ proves the determinant identity.
\end{proof}

At the radial height $z=h$, these expressions simplify to
$D=1$, $a=w$, $b=0$, $c=w'/r$. Thus
\begin{equation}\label{eq:intro-linearization}
 \mathscr L_wu=(wu_r)_r+\frac{w'}r u_{\theta\theta}.
\end{equation}
For $w>0$ it is elliptic where $w'>0$ and hyperbolic where $w'<0$.
A simple zero of $w'$ is a turning circle. Notice that it is the
slope of the profile that changes sign, not the profile itself.

\subsection{Why the endpoint must keep changing type}

It is tempting to take an uncomplicated bump for $w$, solve in its
middle, and flatten its ends. The next observation explains why
that does not preserve a nonzero infinitesimal height.

\begin{proposition}[The monotone-end obstruction]
\label{prop:obs-monotone-tail}
Let $w>0$ on $(r_-,r_+)$ be $C^1$, with $w(r_-)=0$ and $w'\ge0$
near $r_-$. If a continuous endpoint extension of a solution satisfies
\[
 (wf')'-j^2w'f/r=0,\qquad f(r_-)=0,\qquad
 \lim_{r\downarrow r_-}wf'=0,
\]
then $f=0$. The analogous assertion holds at a right endpoint with
$w(r_+)=0$ and $w'\le0$ nearby.
\end{proposition}

\begin{proof}
Integrating the flux equation and putting
$m(r)=\max_{r_-\le t\le r}|f(t)|$ gives
\[
 |f'(r)|\le \frac{j^2}{w(r)}
       \int_{r_-}^r\frac{w'(t)}t|f(t)|\,dt
 \le\frac{j^2}{r_-}m(r).
\]
Integration once more and Gronwall's inequality give $m=0$ on a
collar; ordinary interior uniqueness continues the zero. At the
right endpoint integrate backwards and use
$\int_r^{r_+}|w'|=w(r)$. The proof also applies to complex $f$.
\end{proof}

Consequently a positive profile with only finitely many changes of
monotonicity cannot carry our endpoint-flat nonzero angular mode.
The infinitely many turning circles are forced by this obstruction
within the present construction. Their amplitudes and widths will
decrease fast enough that the profile nevertheless extends smoothly.

\begin{figure}[htbp]
\centering
\begin{tikzpicture}[x=1cm,y=1cm,>=Stealth,font=\small]
\draw[->] (0,0)--(11,0) node[right] {$r$};
\draw[thick,blue!70!black] (0,0)--(.4,0)
 .. controls (.6,0) and (.65,.15) .. (.8,.1)
 .. controls (.95,.05) and (1,.4) .. (1.2,.3)
 .. controls (1.4,.2) and (1.5,.85) .. (1.8,.65)
 .. controls (2.1,.45) and (2.2,1.4) .. (2.8,1.2)
 .. controls (3.5,1) and (3.5,2) .. (4.5,1.75)
 .. controls (5.5,1.4) and (6,2.1) .. (6.8,1.5)
 .. controls (7.5,.9) and (7.5,1.4) .. (8.0,.85)
 .. controls (8.5,.35) and (8.5,.8) .. (8.8,.4)
 .. controls (9.0,.16) and (9.1,.36) .. (9.3,.16)
 .. controls (9.5,0) and (9.6,.12) .. (9.8,0)--(10.5,0);
\node[above] at (5.2,2.15) {positive profile $w$};
\node[below,align=center] at (1.4,-.2) {accumulating turns\\flat extension};
\node[below,align=center] at (8.9,-.2) {accumulating turns\\flat extension};
\draw[->,orange!70!black] (4.8,1.1)--(5.2,1.6);
\node[below] at (5.1,.9) {falling: hyperbolic};
\draw[->,teal!70!black] (3.1,.45)--(3.65,1.5);
\node[below] at (3.1,.4) {rising: elliptic};
\end{tikzpicture}
\caption{Why there are infinitely many interfaces. The drawing shows
only finitely many representative turns, without their actual scales.
On each rising interval the linearized equation is elliptic; on each
falling interval it is hyperbolic. Every derivative still tends to
zero at the accumulation endpoints.}
\label{fig:profile-type}
\end{figure}

\subsection{Two denominators in one finite model}
\label{sec:worked-example}

The next calculation is a guide to the spectral construction, not
an approximation to its full inverse. Set $\epsilon=\sqrt2/100$,
$\rho(X)=\epsilon^2X(1-X)$, and on $[-1,2]$ solve
\begin{equation}\label{eq:example-radius}
 r_{XX}=-\rho r,\qquad r(0)=1,\qquad r_X(0)=1/10.
\end{equation}
Its integral equation bounds $R:=\max|r|$ by
$R\le1.2+8\epsilon^2R$, so $R<1.203$,
$r>0.798$, and $|r_X-.1|\le4\epsilon^2R<.001$.
Hence $w(r(X))=r_X(X)$ is a genuine positive radial profile.
It has two simple turning points, at $X=0,1$.
For an angular mode $j$ and $J=j\epsilon$, its equation becomes
\begin{equation}\label{eq:example-mode}
 L_jU=\frac1{r_X}\bigl(U_{XX}+J^2X(1-X)U\bigr).
\end{equation}

The characteristics have slopes
$d\theta/dX=\pm\epsilon\sqrt{X(1-X)}$ on $0<X<1$.
Their travel angle is
\begin{equation}\label{eq:example-action}
 T=\epsilon\int_0^1\sqrt{X(1-X)}\,dX=\epsilon\pi/8.
\end{equation}
The difference of the two endpoint translations acts by
\begin{equation}\label{eq:example-circle-operator}
 \mathcal D_T f=f(\theta+T)-f(\theta-T),\qquad
 \mathcal D_T\sin(j\theta)=2\sin(jT)\cos(j\theta).
\end{equation}
Thus straightening the two characteristic returns costs the divisor
$\sin(jT)$. For this quadratic irrational travel angle the integer
$2j^2-640000m^2\ne0$, where $m$ is nearest to $j\sqrt2/800$, yields
\begin{equation}\label{eq:example-diophantine-bound}
 |\sin(jT)|\ge (320000j)^{-1}.
\end{equation}
Indeed factor that integer, use
$j\sqrt2/800+|m|\le j$, and then
$|\sin(\pi x)|\ge2\dist(x,\mathbb Z)$.
The inverse of $\mathcal D_T$ therefore loses one angular derivative.

Now impose a radial boundary condition. Under
$x=\sqrt J(X-1/2)$ the homogeneous equation becomes
\begin{equation}\label{eq:example-oscillator}
 (-\partial_x^2+x^2)U=(J/4)U.
\end{equation}
At each end of the finite interval take the solution line obtained
by restricting the corresponding decaying solution of this
whole-line oscillator. These are legitimate boundary lines for
\eqref{eq:example-mode}; the radius itself need not extend outside
$[-1,2]$. The two lines agree exactly at an oscillator eigenvalue.
Since those eigenvalues are $2m+1$
\cite[Section~8.3]{Teschl2009}, the resonance is
\begin{equation}\label{eq:example-spectral-phase}
 J/4=2m+1,
 \qquad\text{or equivalently}\qquad
 jT=(m+\tfrac12)\pi.
\end{equation}
For an inhomogeneous source, multiply by $-r_X/J$, extend by zero
outside $[-3\sqrt J/2,3\sqrt J/2]$, apply the whole-line oscillator
resolvent, and restrict back. This proves that the prescribed boundary
problem has precisely these spectral denominators; it is not just
a comparison of formal phases.

\begin{figure}[htbp]
\centering
\begin{tikzpicture}[x=1.5cm,y=.8cm,>=Stealth,font=\small]
\node[anchor=east] at (-.2,1.5) {characteristic return};
\node[anchor=east] at (-.2,0) {decaying radial lines};
\draw[->] (0,1.5)--(4.4,1.5);
\draw[->] (0,0)--(4.4,0) node[right] {$jT$};
\foreach \x/\t in {0/0,1/\pi,2/2\pi,3/3\pi,4/4\pi}{
 \fill[blue!70!black] (\x,1.5) circle (2pt);
 \node[below] at (\x,-.15) {$\t$};}
\foreach \x in {.5,1.5,2.5,3.5}{
 \fill[orange!80!black] (\x,0) circle (2pt);}
\node[above] at (2,1.7) {$\sin(jT)=0$};
\node[above] at (2,.2) {$jT=(m+\tfrac12)\pi$};
\end{tikzpicture}
\caption{Two distinct resonance sets in the exact finite model.
The horizontal variable is the phase, with marks spaced by $\pi$.
Controlling one row does not control the other. Both must stay
separated when the cells of the infinite construction are joined.}
\label{fig:two-denominators}
\end{figure}

The half-step shift also has a local interpretation. Near $X=0$,
put $y=J^{2/3}X$. Then
\[
 U_{yy}+yU=J^{-2/3}y^2U.
\]
On any fixed $y$-interval the leading solutions are
$\operatorname{Ai}(-y),\operatorname{Bi}(-y)$, with an error
$O(J^{-2/3})$ in value and first derivative for fixed initial data.
Variation of constants against their bounded fundamental matrix
on that interval, followed by Gronwall, gives this estimate.
Their standard asymptotics carry the phase shift $\pi/4$ at each
turn \cite[9.7.5--9.7.12]{DLMFAiry}. The exact oscillator calculation,
rather than this local heuristic, proves
\eqref{eq:example-spectral-phase}.

\subsection{Normalization, conservation and the first residual correction}

The nonlinear unknown is $z=h+qu$. Divide its residual by the small
parameter, using the removable definition
\begin{equation}\label{eq:cart-removable-equation}
 \mathcal A=\mathcal F_q(w,u)=
 \begin{cases}r\mathcal N_h(h+qu)/q,&q\ne0,\\
 \mathscr L_wu,&q=0.
 \end{cases}
\end{equation}
Smoothness at zero follows from the exact formula
\begin{equation}\label{eq:normalized-integral}
 \mathcal F_q(w,u)=\int_0^1 rD_z\mathcal N_h(h+tqu)u\,dt.
\end{equation}
We work with heights even in $\theta$ and invariant under
$\theta\mapsto\theta+\pi$. Nonlinear products preserve this class.
For functions constant at the endpoint collars, Stokes' theorem
applied to the connection gives
\begin{equation}\label{eq:zero-total-integral}
 \int_{r_-}^{r_+}\av{\mathcal F_q(w,u)}\,dr=0,
 \qquad \av f=\frac1{2\pi}\int_0^{2\pi}f\,d\theta.
\end{equation}
Both boundary periods equal $2\pi$. Every differential of this
identity, including metric differentials, has the same compatibility.
Constants in $u$ do not affect the equation.
\label{lem:scalar-integral}

The mean is more transparent at the level of the connection period.
For $U=z-h$ set
\begin{equation}\label{eq:cart-mean-path-flux}
 \mathfrak F_w(U)=\av{\alpha_\theta}-1
 =\left\langle\frac{1+kp_\theta/(1-p^2)}{\sqrt D}\right\rangle-1.
\end{equation}
Its derivative is $r\av{\mathcal N_h(h+U)}$, and its endpoint
value is zero. If $U=Z+m(r)$ with $\av Z=0$, the mean equation
is therefore $\mathfrak F_w(Z+m)=0$. Formula
\eqref{eq:conn-exact-variation} gives
\begin{equation}\label{eq:mean-period-derivative}
 \partial_{m'}\mathfrak F_w(U)=\av a,
 \qquad D_U\mathfrak F_w(0)[U]=w\av{U_r}.
\end{equation}
This linear cancellation is the reason the mean can later gain
twice the decay of the nonconstant height.

The true linearization has one inconvenient feature: away from a
solution, the residual term in
\eqref{eq:cart-determinant-identity} moves the turning circles.
We remove that term while preserving the normal flux.

\begin{lemma}[Fixing the type-change set]\label{lem:type-completion}
On the neighborhood $D\ge1/2$, $Q\ge7w/8>0$, put
\[
 B=E/(rQ),\qquad \bar c=c-Br\mathcal N_h(z),\qquad
 \bar L=\partial_r(a\partial_r+b\partial_\theta)
             +\partial_\theta(b\partial_r+\bar c\partial_\theta).
\]
Then
\begin{equation}\label{eq:cart-completed-normalized-derivative}
 a\bar c-b^2=K_0/D^2,\qquad
 D_u\mathcal F_q\eta=\bar L\eta+
                 q\partial_\theta(B\mathcal A\eta_\theta).
\end{equation}
The corrected operator is symmetric for $dr\,d\theta$, kills
constants, and has the same normal flux as the true derivative.
\end{lemma}

\begin{proof}
The change in determinant is
$aBr\mathcal N_h=ED^{-3/2}\mathcal N_h$, exactly the unwanted
term. The change of operator is
$\partial_\theta(Br\mathcal N_h\partial_\theta)$;
substitute $r\mathcal N_h=q\mathcal A$.
Only the tangential coefficient changed, proving the remaining
assertions.
\end{proof}

The correction vanishes with the equation being solved. Its error
is bilinear in the residual and the increment, the form needed in
Hamilton's implicit theorem. A second correction will fix the
characteristic returns. For reference the present characteristic
slopes on a negative component are
\begin{equation}\label{eq:completed-characteristics}
 \frac{d\theta}{dr}=d\pm v_{\mathrm{char}},\qquad
 d=\frac{E(p-k_\theta)}{rQ},\qquad
 v_{\mathrm{char}}=\frac{\sqrt{-ww'/r}\sqrt D}{Q}.
\end{equation}
At a simple turning circle this speed is a smooth positive factor
times the square root of inward distance. The subsequent coordinate
construction must respect that degeneracy; merely treating all
intervals as uniformly elliptic would lose the geometry of the problem.
